%+PATCH,//WAVEDOC/TEX
%+DECK,wave_by_examples,T=LATEX.
\documentclass[11pt,ngerman,a4paper,twoside]{article}

%\usepackage{epsfig,varioref,titlesec,latexsym,amssymb,amsmath,fancybox,graphicx}
\usepackage{varioref,titlesec,latexsym,amssymb,amsmath,fancybox,graphicx}
%\usepackage{babel}

\usepackage{hyperref}
\hypersetup{
  colorlinks,
  linkcolor=blue,
  filecolor=blue,
  urlcolor=blue,	
  citecolor=blue,
  plainpages=false,
  hypertexnames=false,
}

\usepackage[nottoc]{tocbibind}
\usepackage{soul}
\usepackage{epstopdf}
\epstopdfsetup{update} % only regenerate pdf files when eps file is newer

\usepackage[scaled]{helvet} %Helvetica-Fonts z.B. phvb
\renewcommand*\familydefault{\sfdefault} % sans serif as default
\newfont{\myheadfont}{phvb scaled \magstep 4}

\usepackage{fancyhdr}
\usepackage{lmodern}
\usepackage{amsmath}
\usepackage{mytitlepage}
\usepackage{relsize}

% Die erste Zahl setzt die Schriftgröße, die zweite bedeutet baselinskip
\titleformat{\section}
  {\normalfont\fontsize{11}{15}\bfseries}{\thesection}{1em}{}

\hyphenpenalty=500

%\DeclareMathSizes{10}{8}{8}{7}   % For size 10 text
%\DeclareMathSizes{12}{11}{10}{8}   % For size 10 text
%\DeclareMathSizes{11}{19}{13}{9}   % For size 11 text
%\DeclareMathSizes{12}{20}{14}{10}  % For size 12 text

% Layout Seitenkopf
%\addtolength{\headwidth}{\marginparsep}
%\addtolength{\headwidth}{\marginparwidth}

\setcounter{secnumdepth}{3}
\setcounter{tocdepth}{3}

\newfont{\bss}{bessy scaled \magstep 4}

\setlength{\unitlength}{1cm}

\clubpenalty=10000
\widowpenalty = 10000
\displaywidowpenalty = 10000

\textwidth 16. cm
\topmargin .0 cm
\textheight 22.5 cm
\parskip 0.3 cm
\sloppy %Gegenteil von \fussy
\oddsidemargin 0.0 cm
\evensidemargin 0.0 cm
\parindent=0.cm

\usepackage{xcolor}
\definecolor{orange}{rgb}{1.,0.8,0.4} % rgb
\definecolor{hellgrau}{gray}{0.9}
\definecolor{schwachgrau}{gray}{0.95}
\definecolor{grau}{gray}{0.4}
\definecolor{dunkelgrau}{gray}{0.2}

\begin{document}

\thispagestyle{empty}

\begin{center}
\LARGE \textbf{WAVE by Examples} \\[1cm]

\textbf{ \Large Michael Scheer} \\[0.cm]
\textbf{ \Large Helmholtz-Zentrum Berlin} \\[0.75cm]

\textbf{ \Large \today }\\
\end{center}

\newpage
\thispagestyle{empty}
{~}

%\newpage

\setcounter{page}{1}
\tableofcontents

%\newpage

\pagestyle{fancy}

%\lhead[\bfseries\thepage]{\bfseries\rightmark}
\lhead[\textit \thepage]{\textit \rightmark}
\chead[]{}
%\rhead[\bfseries\leftmark]{\bfseries\thepage}
\rhead[\textit \leftmark]{\textit \thepage}

\lfoot[]{}
\cfoot[
Version 0.00/02
]{}
\rfoot[]{}

%\large

\include{hyphenation}
\include{abbrevs}

%+PATCH,//WAVE/TEX
%+DECK,introduction,T=LATEX.

\renewcommand*\thesection{\Roman{section}}
\section{Introduction}

\W~has been developed at BESSY over the years to track electrons
through magnetic fields and to calculate spontanous synchrotron radiation
\cite{dok}.
This document illustrates the usage of \W~and makes you
familiar with it by following the given examples.

\W is mainly controlled by an ASCII input file named \wie.
For each example, there is a script \gtis{wave\_example\_n.sh} to
create \wi from \wee, where \gtis{n}~is the number of the example. Once
\wi is created by the corresponding script, you may compare
\we and \wi to see, which variables have changed.

There is a graphical user interface (GUI), a Python script called
\WSH to setup \wie, to run \W, and to visualize the
results. This is explained in the first example. However, not all
options for \W are available in the GUI.

Once \W is installed, you should have a \W directory and shell
variable \W pointing to it. It is assumed, that the user knows to edit
a file and how to open a terminal, e.g. \gtis{bash} under Linux  or powershell and under
Windows. As an alternative one may run the Python shell \gtise{idle}.

WAVE and the related software is published under the GNU General Public
License. To get WAVE, please contact the author or get it from github:

~~~~~\gtis{git clone https://github.com/MicScheer/wave.git wave}


\section{Getting Started}

The examples in this document refer to a \gtis{LINUX} machine, however, they
should be easy to understand and adapted to a \gtis{WINDOWS} computer.

The main differences are, that the slashes in the file and path names must
be replaced by backslashes, shell scripts ($*.sh$) must be replaced
by  $*.bat$ files, and the system variables $\$V$ must be
replaced by $\%V\%$.

\subsection{Running WAVE}

Under \gtis{WINDOWS} open e.g. a DOS command shell (cmd.exe), goto to
the directory where \W is installed and excute
\gtis{bat\textbackslash set\_wave\_environment.bat}\bsp. Then enter
\gtis{wave} or \gtise{waveshop}. But as mentioned before, we now refer to \gtise{LINUX}.

First define a shell variable \dw~ containing the path to your
\W-directory:

~~~~~\gtise{export WAVE=YOUR-WAVE-DIRECTORY}

~~~~~\gtise{export WAVE\_INCL=\dw}

It is recommended to put these lines in a script, see e.g.

~~~~~\senv

and to excute it at the beginning of a \W-session.

To run \W the traditional way is to is to enter in a terminal

~~~~~\gtis{cd \dw/bin/stage}

~~~~~\gtis{../bin/wave.exe}\\

It is recommended to use the script

~~~~~\gtis{\dw/stage/wave}

to run \W.

WAVE should run and write some output to your terminal (Fig.~\ref{gsfig3}).

\begin{figure}[htb]
    \centering
    \includegraphics[width=0.75\textwidth]{pics/wave_run.png}
    \caption{Terminal output of \textit{WAVE}}
    \label{gsfig3}
\end{figure}

The program writes some output to the screen and the output text files
\wo and
\mhbe. The file \wo contains the most important
in- and output of the run, while \mhb contains the histograms and
n-tuples to be plotted. Please open these files and have a look{\ldots}

If the program crashes, it must probably be recompiled in your
environment by

~~~~~\gtis{python3 \dw/python/make\_wave.py}

if you have \gtis{gfortran} installed, or by

~~~~~\gtis{python3 \dw/python/make\_wave\_intel.py}

if you have the Intel compiler \gtis{ifx} installed.

If these Python scripts are run with an argument, they run in a verbose mode.

\subsection{Plotting with Python}

A convenient way to visualize the histograms and
n-tuples is to use a Python scripts

~~~~~\gtis{python3 -i ../python/waveplot.py}

First, we check the installation of Python. Enter:

~~~~~\gtis{python3}

If you get a \gtis{Python} command prompt, \gtis{Python3} is
installed. Note: Under \gtis{WINDOWS} the command is probably
\gtis{python} instead of \gtis{python3}. However, the version $3.x$ of {\it
Python} is required. You can leave the \gtis{Python} shell with:

~~~~~\gtis{quit()}

If the test was successful, try the same with ipython3:

~~~~~\gtis{ipython3}

If there is no command \gtis{ipython3}, you can install it (given pip3
is installed):

~~~~~\gtis{pip3 install ipython3}

\subsection{The Graphical User Interface \WSH}

There is a graphical user interfaces (GUI) based on \gtis{Python} and
\gtis{Tkinter}.

To run the GUI \WSH (Fig.~\ref{gsfig1}), you need \gtis{python3} or
\gtie{ipython3}. In addition,
the file \gtis{waves.wvs} must be present in the
working directory. The \gtis{Python} script \gtis{../python/waveshop.py} is started by
the command alias \gtis{waveshop}. Thus enter

~~~~~\WSH

or, if the alias is not defined accordingly, by

~~~~~\gtis{python3 -i ../python/waveshop.py}

\begin{figure}[t]
    \centering
    \includegraphics[width=0.95\textwidth]{pics/waveshop_gsfig1.png}
    \caption{Main window of \WSH. The upper menu items
    are used to set-up \wi and to run \W, the lower buttons to plot the
    results.}
    \label{gsfig1}
\end{figure}

%\newpage
\subsection{The Input File \gti{wave.in}}

The GUI for \W covers many standard calculations, however,
to get the most out of \W, it is necessary to study and
understand the input file \wie.

\begin{figure}[htb]
    \centering
    \includegraphics[width=0.7\textwidth]{pics/ex1_wave_in.png}
    \caption{First part of the input file \textit{wave.in}}
    \label{gsfig6}
\end{figure}

We open \wi with an editor and search \contrl. This keyword
belongs to a \gtis{FORTRAN} namelist. A namelist begins with such a keyword
and ends with \dend. The file \wi is actually nothing than a list of
namelists.

A namelist is a list of variables, read by the program. The position
of the variable in the namelist does'nt matter. Missing variables are
initiliazed as zero or an empty string. For the name of the variables
the \gtis{FORTRAN} conventions aplly, i.e. variables beginning with \gtis{I, J, K,
L, M, N} are integer, all others are strings, floating point, or complexe
numbers. The integer must not have a decimal point, while the floating
point  and complexe numbers onces must have. The same
variable may occur several times within a namelist. Then the last
value is accepted. If a whole namelist occurs several times, only the
first one is read. Comments must begin with~"!".

The namelist \contrl is the most important one of \wie. It
controlls the program flow, the tracking, the field selection and so on.

\newpage
A short example: Look for a line, containing

\cboxs{KELLIP=1}

This variables activates an analytical model of an elliptical
undulator. The comment in this line mentions the namelist \scap{\$ELLIPN}.
Hence, we search \scap{\$ELLIPN}.

\begin{figure}[htb]
    \centering
    \includegraphics[width=0.7\textwidth]{pics/ellipn.png}
    \caption{The namelist \gtis{ELLIPN} of
    the elliptcal undulator model}
    \label{gsfig7}
\end{figure}

\newpage
Here you could change the parameters of the selected elliptical
undulator.

\ul{Note}: The GUI \WSH can't handle multiple occurances of
variables, i.e. it always takes the first value. So it's better to
avoid it. If you want to keep a copy of an old value, change the line
into a comment line.

Editing input files might appear a little old-fashioned today, but has
the big advantage of easy documentation and allows to run and control
the program from scripts, which can be very effective e.g in
optimization procedures.

\renewcommand*\thesection{\arabic{section}. \hspace{-4mm}}
\setcounter{section}{0}

\section{Example: Selecting a Magnetic Field Model}

%\vspace{-1cm}

After

~~~~~\senv

our starting point is \wee. We copy it to \wie:

~~~~~\gtis{cp \$WAVE/input/wave.example wave.in}

and enter

~~~~~\WSH

The GUI pops up, select in the upper menu bar \scap{Comment} and
to change the user title

~~~~~\scap{WAVE.EXAMPLE}

to

~~~~~\gtis{Planar Undulator}

Hit \scap{OK} to leave the widget. Select \gtise{Sel. Mag.}.

A menu pops up:

\begin{itemize} \renewcommand\labelitemi{--}
\item
Click

\witem{Elliptical undulator (analytical model)}

to deselect the current model and

\witem{Halbach undulator/wiggler with end poles}

to select the planar undulator model.
\end{itemize}

\begin{figure}[htb]
    \centering
    \includegraphics[width=0.6\textwidth]
    {pics/halbasy_undulator.png}
    \caption{Parameters of the undulator}
    \label{ex1fig1}
\end{figure}

Now we set up an undulator model with \gtis{99} main poles and two
endpoles, i.e. a \gtis{50} periods undulator. The period length is
\gtise{0.033m}.
In addition, the third harmonic should be at \gtis{500eV}
(Fig.\ref{ex1fig1}).

\underline{Note:}
\begin{itemize}
\item
\vspace*{-0.5cm} Selecting a harmonic overrides the magnetic field value given in the shown
menu.
\item
For this model, \gtis{z} is the longitudinal
coordinate, while the WAVE coordinate system is such, that \gtis{x} is the
longitudinale axis.
\end{itemize}

We close the last two menus and go back the main input menu.

Next, we adjust the observation pinhole.

\begin{itemize} \renewcommand\labelitemi{--}

\item
Click

\witem{Spectrum calculations}
\witem{Pinhole definition}

and set the pinhole width and height to \gtis{0.002}

\item
Hit \gtis{OK} and choose:

\witem{Spectral range}

Set the lowest and highest photon energies to \gtis{495.0} and
\gtis{505.0} eV, respectively.

\end{itemize}

We run \W by clicking \gtis{Run} in the toolbar.

\begin{figure}[htb]
    \centering
    \includegraphics[width=0.9\textwidth]{pics/ex1fig3_overview.png}
    \caption{Overview of plots}
    \label{ex1fig3}
\end{figure}

In the terminal we see some warnings, which can be ignored. They could
be relevant in the case of very precise calculations. To avoid them,
one could increase the number of grid points in the Pinhole menu from 21 to
51, but that slows down everything for nearly no benefit.

Click the button \gtis{Plot}, select \gtis{Overview}, and check the results.

Before we finish this examples, we open the output file \wo~with a
text editor.

\begin{figure}[htb]
    \centering
    \includegraphics[width=0.9\textwidth]{pics/ex1_wave_out.png}
    \caption{The output file \textit{wave.out}}
    \label{ex1fig2}
\end{figure}

This file contains a summary of the input variables and results of the
calculations, as well as warnings and error messages. At the top you
see your user comment (the variable \scap{CODE} in \wi) and the current
run number, which is stored in the file \scap{WAVE\_CODE.DAT}. The
output file is a good place to check the input parameters of the
magnetic field model, the beam parameters etc..

As a beginner, you
should at least once study the file thorougly, although it might
contain some information you don't understand at this moment{\ldots}

%\newpage
Now it's up to you, to fumble around with the parameters of the
menus described above and watch the effects. In case of trouble, you
can always go back by copying
\gtis{\$WAVE/input/wave\_example\_1.in}
to
\gtis{wave.in}.

Proposed exercise: Compare \wi and \we.
 % Selecting a Magnetic Field Model
\section{Example: Tracking through Magnetic Fields}

This examples illustrates the tracking mechanism of \W (and by the way
the use of the command line{\ldots} for the plotting.) After

~~~~~\senv

We start from
\we and use the python script
\gtis{../python/wave\_example\_2.py}
to create \wie. But before, we have a look on it. We open
\gtis{../python/wave\_example\_2.py} with an text editor (Fig. \ref{ex2fig1}):

\begin{figure}[htb]
    \centering
    \includegraphics[width=0.9\textwidth]
    {pics/wave_example_2_py.png}
    \caption{First part of the script to set-up \textit{wave.in} for the second example
    \textit{wave.out}}
    \label{ex2fig1}
\end{figure}

The list \gtis{ReplList} contains the strings to be replaced in \wie. Each
line contains two strings. The first string
in \wi is replaced by the second.
Most of the variables you known from example 1,
but the variable
\ispec belonging to the namelist \contrl is new.
Changing \ispecone to \ispecnull, switches the spectrum
calculations off. Also new is the variable \gtis{BYGOFF}, which is a
global field offset added to the undulator field, e.g. to simulate the
earth field.

Now we execute the script,
copy the resulting file, and run \W:\\

\cboxs{
python3 ../python/wave\_example\_2.py \\
cp ../input/wave\_example\_2.in wave.in \\
wave\\
}

On the terminal, two warnings appear, indicating that the trajectory
is not correctly adjusted. This means, a final kick and a final offset
remain. We start the plotting GUI:

\cboxs{wplot\\}

Click the button \gtis{Plot}, select \gtis{Overview}, and check the results.

We see in the overview Fig. {\ref{ex2fig2}} that
the trajectory is distorted as expected due to the external field,

\begin{figure}[htb]
    \centering
    %\framebox(12,9){
      \includegraphics[width=0.7\textwidth]
      {pics/trajectory_ex2.png}
    %}
    \caption{The trajectory with \gtis{XINTER=-9999.}}
    \label{ex2fig2}
\end{figure}

Sometimes, it's useful to adjust the trajectory, e.g. if you want to have
a symmetric one with respect to the center of the straight section or
if for a given field model the correct endpole constellation is
missing. This can be achieved with by tuning the electron parameters in the
namelist \contrle:

The most interesting variable is \gtis{XINTER}, which defines the
x-coordinate the other variables like \ystart or \vzin refer to.

Edit \wie and set

\cboxs{XINTER=0.}

Then we save the file and enter:

\cboxs{wave \\
wplot \\
nz()
}\\

and see (Fig. \ref{ex2fig4}) a trajectory, where the electron passes the origin with zero
slope according to:

\begin{figure}[htb]
    \centering
    \includegraphics[width=0.8\textwidth]
    {pics/ex2_xinter.png}
    \caption{Electron variables}
    \label{ex2fig3}
\end{figure}

\cboxss{
\bsp ~YSTART=0.\\
ZSTART=0.\\
VXIN=1.\\
VYIN=0.\\
VZIN=0.
}

at

\cboxss{XINTER=0.}

Thus, if \scap{XINTER} is not \gtis{-9999.}, WAVE tracks
first the electron from \scap{X = XINTER} back to
\xstart and then starts the actual tracking through the magnetic
field.

You can save the picture with
the corresponding toolbar button or typing ''s'' when the cursor is
over the plotting window or by typing e.g.

\cboxs{pplot("trajectory.png")}\\
or

\cboxs{pplot("trajectory.pdf")}

in the Python shell. (If the overview is plotted a file \scap{
WaveOverview.pdf} is written automatically.) The last plotting command \\

 \cboxs{nz()}

can of course also be triggered from the GUI menu. The ''n'' referes
to the fact, that the data stem from an n-tuple. You can also try \\

 \cboxs{nby()}

and so on. You might have alread noticed, that the plotting commands
executed by the GUI are echoed in the Python shell.
The command echoes show you
the syntax and parameters of the plotting commands applied. So you can
learn short-cuts for plotting from the Python shell.

If you don't see these command echos,
the configuration file \gtis{waveplot.cfg} is missing in your working
directory. You find an example in the input-directory \gtis{\$WAVE/input}.

If \gtis{waveplot.cfg} is present, plotting options are read from this
file. E.g. you can set the window size, trigger
saving of all plots or writing plotted
data to ASCII-Files. To see all options, check this file.

Some options can also be toggled in the GUI or by commands, e.g.

 \cboxs{optecho(1)}

 or

 \cboxs{optecho(0)}

 or

 \cboxs{optstat(1)}

 or

 \cboxs{optstat(0).}


\begin{figure}[htb]
    \centering
    %\framebox(12,9){
    \includegraphics[width=0.7\textwidth]
    {pics/ex2_xinter_0.pdf}
    %}
    \caption{Adjusted trajectory with \gtis{XINTER=0.}.}
    \label{ex2fig4}
\end{figure}

Note: If you calculate synchrotron radiation from this trajectory,
i.e. \ispecone, the
back tracking at the beginning does, of course, not contribute to the
radiation or irradiated power of the device. By the way, this power is
given in the output
file \woe. The power is normalized to the beam current
\dmycurr in the namelist \contrl. In \wo you also find the energy loss
of the electron having passed the magnetic field. This is controlled by
\ieneloss also namelist \contrl. For very long undulators or
wigglers, the energy loss influcences the trajectory and the spectrum
of synchrotron radiation. For a single, typical insertion devices,
this can be neglected, i.e. \ienelossnull can be used.
 % Tracking through Magnetic Fields
\section{Example: The Concept of Source Points \label{ex3}}

Synchroton radiation is mainly emitted into a narrow cone on the scale
of $1/\gamma$. Consider a wave-length shifter (WLS) with a main pole and two
endpoles. Such a device has a deflection angle that is usually much
larger than $1/\gamma$. Therefore, an observer at point $P$ sees up to
three sources of radiation from the WLS, namely the points of the
trajectory, where the electron points to $P$ as sketched in Fig.
\ref{ex3fig1}. These points, we call source points. Actually, they are
not really points, but pieces of the trajectory due to the finite
opening angle of the radiation cone.

\begin{figure}[htb]
    \centering
    \includegraphics[width=0.8\textwidth]
    {/home/scheer/wav/doc/pictures/gut-spect-1a.pdf}
    \caption{Source points of a wave-length shifter for
    observation point P}
    \label{ex3fig1}
\end{figure}

Let's have a closer look. Enter in the terminal:

\cboxs{
. \$WAVE/shell/set\_wave\_environment.sh\\
cp ../input/wave.example wave.in\\
waveshop
}\\

\begin{itemize} \renewcommand\labelitemi{--}
\item
The GUI pops up, select in the upper menu bar \gti{Comment} and
change the user title

\vitemt{User title}{
Example 3: The Concept of Source Points
}

\item
Click \gtis{Track} in the toolbar

and set
\vitem{number of steps per meter}{100000}

to increase the precision of the spectrum calculations in the
undulator mode. Hit \scap{OK}.\\

\item
Click \gtis{Sel. Mag.} in the toolbar and in the menu

\witem{Elliptical undulator}

\witem{Halbach undulator/wiggler with end poles}

{~}\\[0.5cm]
and set the peak field to \gtis{6.0T}, the number of main poles to \gtis{1}, the
period length \gtis{0.7m}. Hit \scap{OK}, again \scap{OK}.\\

\item
Click \gtis{Spectra} in the toolbar and in the menu

\witem{Pinhole definition}
\witem{single point}

{~}\\[0.5cm]
Click \scap{OK}.\\

\item
Click
\witem{Spectral range}

{~}\\[0.5cm]
and set the lowest photon enery to \gtis{50.}, the highest to \gtis{25000.}, and
the number of photon energies to \gtis{2001},  then hit \scap{OK}
(Fig. \ref{ex3range}).

\begin{center}
\begin{figure}[htb]
    \centering
    \includegraphics[width=0.7\textwidth]
    {/home/scheer/wav/doc/pipy/wave_example_3_range_2.png}
    \caption{Setting of the spectral range}
    \label{ex3range}
\end{figure}
\end{center}

\end{itemize}

%\newpage
Close the menu and click \gtis{Run} in the toolbar and plot the overview.\\

\begin{figure}[htb]
    \centering
    %\framebox(12,9){
      \includegraphics[width=0.7\textwidth]
      {/home/scheer/wdp/ex3_wls_undulatormode_overview.pdf}
    %}
    \caption{The spectrum of an WLS calculated in the
    undulator mode}
    \label{ex3fig2}
\end{figure}

%{~}\\[1cm]
The spectrum (right plot in the overview Fig. \ref{ex3fig2}) looks a little
weird. We leave the GUI, edit \wi and set

\cboxss{FREQLOW=2000.\\FREQHIG=2010.\\NINTFREQ=201\\IEFOLD=3}

To run \W, enter:

~~~~~\gti{wave}

%{~}\\[-0.5cm]

When \W has finished,  enter

~~~~~\gti{wplot}

%\newpage

We choose in the plotting GUI

\cboxs{Plot/Spectra/Flux-density/Fd}\\

The spectrum (red) in Fig. \ref{ex3fig3a}
is not as smooth as we would expect from a wave-length shifter.
The synchrotron radiation of a single electron is fully
coherent, thus there are interferences between the radition of the different poles of
the WLS. But the
phase advance between the poles is so high, that the
corresponding huge number of oszillations can not be resolved by the
integration algorithm of \W. The observed spectrum
is therefore a mixture
of real interferences and a numerical noise depending on the
integration parameters.

However, the
energy-spread of a real beam smears out the spectrum.
To plot the smeared spectrum goto to the plot GUI and select:

\cboxs{Options/Same Zone\\
Options/Next Color\\
Menu/Spectra/Flux-density/Fd with e-spread}

The resulting spectrum (blue line) is smooth as would be expected from
a wave-length shifter.

\begin{figure}[htb]
    \centering
    %\framebox(12,9){
      \includegraphics[width=0.7\textwidth]
      {/home/scheer/wav/doc/pipy/ex3_wls_undulatormode_a.png}
    %}
    \caption{\parbox[t]{10cm}{The spectrum of an WLS calculated in the
    undulator mode. Red: Without e-spread, blue: with e-spread.}}
    \label{ex3fig3a}
\end{figure}

Terminate the plotting and start the GUI \WSHE.

In the toolbar hit \gtis{Spectra} and select in the menu "Spectrum calculations"

\witem{Main modes} \witem{wiggler mode 2}

Hit \scap{OK} and run \W.

To plot the spectrum including the energy-spread, we choose in the plotting
GUI

\cboxs{Menu/Spectra/Flux-density/Fd with e-spread}

or enter in the Python shell

\cboxs{hcfluxden("fde")}

\begin{figure}[th]
    \centering
    %\framebox(12,9){
      \includegraphics[width=0.7\textwidth]
      {/home/scheer/wav/doc/pipy/wave_example_3_wiggler.png}
    %}
    \caption{The spectrum of a WLS calculated in the wiggler
    mode including the energy-spread of the beam}
    \label{ex3fig5}
\end{figure}

Fig. \ref{ex3fig5} shows the same result, as in the previous run with the
undulator mode. We conclude, that the interferences between the poles
of the WLS are real, but can not be calculated correctly and can not be
observed for real beams with a finite energy-spread.
Hence, it is justified to neglect them, i.e. to calculate the radiation
from each pole seperately and to sum up the contribution of all poles
incoherently. This is what the wiggler mode of \W does. The decision,
if the wiggler mode should be applied (and especially together with
Schwinger's formular) depends on the  regime, where the spectrum is
calculated. For a wiggler, this  means the undulator mode is needed in
the spectral range of the lower harmonics. In doubt, try the different
modes and have a look on the results.

\begin{figure}[th]
    \centering
    %\framebox(12,9){
      \includegraphics[width=0.7\textwidth]
      {/home/scheer/wav/doc/pp/wave_example_3_sources.png}
    %}
    \caption{Source points of a WLS}
    \label{ex3fig6}
\end{figure}

In the plotting GUI, choose:

\cboxs{Menu / Fields and Trajectory / Sources / z(x)}

As shown in Fig. \ref{ex3fig6}, \W has found three source points. How
does it work? \W tracks first the electron and stores the trajectory.
Then it checks for all steps of the trajectory, if a beam with a given
opening angle passes at least partically a collimator.
The first step satisfying this, is marked as the beginning of a
source, the last as the end.

%\newpage

This procedure and the set-up of the collimator are controlled by the variables of the namelist
\scap{\$COLLIN}. The main modes of \W namely the undulator and
wiggler modes set default values for these variables, which
should work in most cases. In addition, the main modes control
the intergration method \scap{ISPECMODE} and the number of integration steps
\scap{NLPOI} used to calculate the
spectrum.

\underline{Proposed exercise:}

Try the expert mode, i.e.
\scap{IUNDULATOR=0, IWIGGLER=0}
and play with the variables
\scap{WGWINFC, NLPOI} and the variables of the namelist
\scap{\$COLLIN}.
 % The Concept of Source Points
\section{Example: APPLE-II Undulator} \label{ex4}

The most realistic undulator model
\footnote{See also example \ref{ex11} ~ for an even better model by interfacing
to \Ue. \U allows to take into account the interaction
between the neighbouring magnets, while WAVE only consideres static
magnets with $\mu = 1$. This interaction might affect especially the
first and second integrals, i.e. the kick and off set at the end of
the trajectory.}
of \W is an undulator build of
permanent magnet blocks (rare earth componds like NdFeB, shortly REC).
The field of these blocks is calculated with the current sheet method.
Due to the large number of blocks of an undulator the field
calculations are rather slow, but the field is very realistic, and an
undulator can be modelled with all its features like gap and shift
changes.

The namelist \scap{\$REC} holds the variables of the undulator. It also
includes gap taper and field errors. To set up the input file \wie, enter:

\cboxs{
\senv\\ {\grau
python3 ../python/wave\_example\_4.py\\
cp \dw/input/wave\_example\_4.in wave.in}
}

After the \W run, you'll find a file named \gtis{rec\_plotm.eps}, which
shows the set-up of the undulator (Fig. \ref{ex4fig1}). \\[-2cm]

\begin{figure}[ht]
    \centering
    \includegraphics[width=0.8\textwidth]
    {/home/scheer/wav/doc/pictures/wave_example_4_recplotm_1.pdf}
    \vspace{-1.5cm}
    \caption{REC model of an undulator}
    \label{ex4fig1}
\end{figure}

The figur shows the upper and lower rows of magnets as top views and a
side view of both. To see details we
click \gtis{Conf. Mag.} in the toolbar of \WSH and \\

\witem{REC-undulator (Option KBREC)}
{~}\\[-1.5cm]
\witem{Options}
{~}\\[-1.5cm]
\witem{Plotting options}

and set the plotting range

\cboxss{RPLXMN = -920.\\
RPLXMX = -850.
}

and rerun \W.\\

\begin{figure}[ht]
{~}\\[-2.5cm]
    \centering
    \includegraphics[width=0.8\textwidth]
    {/home/scheer/wav/doc/pictures/wave_example_4_recplotm_2.pdf}
    \vspace{-1.5cm}
    \caption{Endpoles of the REC model}
    \label{ex4fig2}
\end{figure}

Now, you see the details of the endpoles, the shift, and the
magnetization of the blocks. More complex situtation can be handle by
an input file containing arrays of magnet. This will be explained in
one of the following examples.

The overview of this run is shown in Fig. \ref{ex4fig3}. It shows the
elliptical trajectory and the corresponding first harmonic.\\

\begin{figure}[ht]
    \centering
    %\framebox(11.5,8){
      \includegraphics[width=0.7\textwidth]
      {/home/scheer/wav/doc/pp/wave_example_4_overview.pdf}
    %}
    \caption{Example of an APPLE-II undulator}
    \label{ex4fig3}
\end{figure}

\newpage
\underline{Proposed exercise:}

In the option menu of the REC-Undulators is a menu to define errors of
the magnetic blocks. Give it a try, and check also the magnetic field
errors option in the menu of the global magnetic field options.
 % An APPLE-II Undulator
\section{Example: Simple Build-In Optimization} \label{ex5}

{~}\\[-0.5cm]
A very common situation is, that you want to tune input parameters
with respect to a figure of merit. \W has build in two simple
procedures to accomplish that. The first is an iterativ estimation
for the case of one parameter to adjust, the second a Monte-Carlo for more
than one parameter.

Consider example \ref{ex4} For the given gap, we want to find the
shift such, that the vertical and horizontal field of the undulator are
the same, i.e. the device acts as an circular undulator.

In a terminal, we enter:\\

\cboxs{
\senv\\ \grau
python3 ../python/wave\_example\_5.py\\
cp ../input/wave\_example\_5.in wave.in\\
wave\\
}

We see in \woe, that the max. vertical and horizontal fields \scap{BYMAX} and
\scap{BZMAX} differ. Now, we
adjust the shift parameter of the undulator. Enter:\\

\cboxs{cp ../input/wave\_example\_5.adjust wave.adjust\\
cp ../python/wave\_adjust\_example\_5.py wave\_adjust.py\\
}

Have a look on \gtis{wave.adjust} (Fig.\ref{ex5fig2}).
When a file like this is passed as an argument to \We, \W will adjust the
parameter \scap{URSHIFT(1)} of \wi by a Monte-Carlo search. Of course,
this also works with other filenames and variables of \wie.  During
the optimization
procedure \W executes the script \gtis{wave\_adjust.sh}, which calls
another instance of \W and \gtise{wave\_adjust.py}.\\

\begin{figure}[htb]
    \centering
    \includegraphics[width=0.8\textwidth]
    {/home/scheer/wav/doc/pp/wave_example_5_wave_adjust.png}
    \caption{Input file for the optimization procedure}
    \label{ex5fig1}
\end{figure}

After the run of this instance,
the command given in the last line of \gtis{wave.adjust} is executed by
the master process.
In our example, this is the python script  \gtis{wave\_adjust.py}.
What does this script (Fig.\ref{ex5fig2}) do? \\

\begin{figure}[htb]
    \centering
    \includegraphics[width=0.8\textwidth]
    {/home/scheer/wav/doc/pp/wave_example_5_wave_adjust_py.png}
    \caption{Example of a program to calculate the figure of
    merit}
    \label{ex5fig2}
\end{figure}

It opens \woe, and looks for the lines containing the keywords like
\scap{BYmin}. When a matching line is found,
the program reads from the line the values of  \scap{By} or \scap{Bz},
respectively,
and calculates the \gtis{$\chi^2$} to be minimized. The resulting \gtis{$\chi^2$}
is written to \gtis{wave\_adjust.dat} and evaluated by the
master process.

Let's try. Enter:

\cboxs{cp wave.in wave\_adjust.in\\
wave wave.adjust \\
}


\W loops and writes \gtis{wave\_adjust.out}.
The data in this file show a minimum for the shift of about \gtis{19mm}
(Fig.\ref{ex5fig3}). This minimum value found is also used to re-run
WAVE.

%\newpage
Start the plotting GUI by entering:

\cboxs{wplot}

and enter in the python shell:\\

\cboxs{adj = ncread("adj","sh:chi2","wave\_adjust.out")\\
nplot(adj,"sh:chi2")\\
txyz("wave\_adjust.out","Shift[mm]","Chi2")}

\begin{figure}[htb]
    \centering
    %\framebox(12,9){
    \includegraphics[width=0.7\textwidth]
    {/home/scheer/wav/doc/pp/wave_adjust_mc.pdf}
    %}
    \caption{Random variation of URSHIFT(1)}
    \label{ex5fig3}
\end{figure}

The plot you get, might differ from that shown here, due to different
seeds of the Monte-Carlo search.

To get a precise value,
we edit \gtis{wave.adjust}, and set the number of parameters to \gtis{-1}
and the boundaries to \gtis{18.0} and \gtis{20.0}.
Then we repeat

\cboxs{cp wave.in wave\_adjust.in\\
wave wave.adjust\\
}

and find now in \wi the best value of about \\

\scap{URSHIFT(1)= 18.6}\\

Since we have required a precision for the minimum
\gtis{$\chi^2 \le 10^{-6}$}, the ratio of the vertical and horizontal field
now agrees on the scale of \gtis{$\sqrt{10^{-6}}$}.

The effort in this example might appear quite high for such a simple
probleme, which could be solve sufficently by a quick interpolation.
However, it shows the usage of scripts and the advantage of command
line driven optimization procedures.
 % Build-in adjustments and optimizations
\section{Example: ''Real'' Beams, Part I} \label{ex6}

Up to now, we have only considered single electron beams, except in example
\ref{ex3}, where we have seen the effect of the energy-spread of the
beam. In this example, we learn about the first of two methods how \W~takes the
effects of real beams, i.e. emittance and enery-spread into account.\\

In modern low emittance storage rings the spatial and spectral
distribution of the synchrotron radiation does not vary very much for
the different electrons in the beam phase-space. Thus, the shape of
the radiation cone is more or less the same, it is only shifted
according to the position and slope of the considered electron.
Therefore, it is possible to take the beam emittance into account by
folding the radiation cone of the reference electron with
a distribution that reflects the phase-space. This distribution is
usually a 2D-Gaussian with
\begin{eqnarray*}
\Sigma_y^2 &=& \sigma_y^2 + d^2 \times \sigma_y'^2   \\
\Sigma_z^2 &=& \sigma_x^2 + d^2 \times \sigma_x'^2
\end{eqnarray*}

where $\sigma_y,\, \sigma_y', \sigma_z, \, \sigma_z'$ denote the
vertical and horizontal sizes
and divergences of the beam  and $d$ the distance between the source and the
observation plane. Enter:\\

\cboxs{\senv\\ \grau
python3 ../python/wave\_example\_6.py\\
cp ../input/wave\_example\_6.in wave.in\\
wave\\
wplot\\
}

Use the plotting GUI to produce the Fig. \ref{ex5fig4} or enter in the
python shell:

\cboxs{zone(1,2,1)\\
ndistpinh()\\
nextzone()\\
ndistpinv()
}

If you want to change parameters for varies runs, but
plot always the same, you can put these lines into a file

\gtis{waveplot\_startup.py}

The commands in this file are exceuted at the start of \wple, until the
end of the file is reached, or a line with a key-word \gtis{EOF} is
found.
Of course, this can also handle the plotting zones by using the \gtis{Options} menu of
\wple.

\begin{figure}[htb]
    \centering
    %\framebox(10,8){
          \includegraphics[width=0.6\textwidth]
	  {/home/scheer/wav/doc/pp/wave_example_6_hvcut.pdf}
    %} % end framebox
    \caption{Horizontal and vertical cuts through the
    flux-density distribution}
    \label{ex5fig4}
\end{figure}

The mean value of the histgrams agree with the position of the
observation pinhole as given in \woe. We see, that the position of the
pinhole is shifted. This has been achieved by setting\\

\cboxss{YSTART=56.1e-6\\
ZSTART=648.e-6\\
VYIN=12.87e-6\\
VZIN=74.4e-6\\
PINCEN(2)=-9999.\\
PINCEN(3)=-9999.\\
}

in \wie . The values for the start conditons of the electron correspond
to three sigmas of the beam sizes and divergences as given in the
namelist \scap{\$WFOLDN}. The values for the center of the pinhole
trigger the shift of the pinhole, as is explained in the namelist
\scap{\$PINHOLE}.

%\newpage
When we set in \wi\\

\cboxss{YSTART=0.0\\
YSTART=0.0\\
ZSTART=0.0\\
VYIN=0.0\\
VZIN=0.0\\
}

as usual, we get nearly the same histograms, except that they centered
on-axis. Hence, the ansatz of the folding procedure is verified.

A similar folding procedure is applied for the spectrum to take into
account the energy-spread of the beam (\scap{ESPREAD} in the namelist
\scap{\$WFOLDN}). This procedure makes use of the fact, that the
harmonics of a undulator depend on the square of the beam energy
$E_{beam}$. Thus, the shift of the spectrum is in first order given
by\\[-1.5cm]

\begin{eqnarray*}
\frac{\Delta E_{phot}}{E_{phot}}  &=&  2 \times \frac{\Delta E_{beam}}{E_{beam}}
\end{eqnarray*}
while the shape stays in a good approximation the same.

\underline{Proposed exercise:}\\

Try to verify the last statement by using \W.
 % Real Beams, Part 1
\section{Example: Global Magnetic Field Options} \label{ex7}

\W~offers in the namelist \scap{\$B0SCBLOBN} some options,
e.g. shifting, scaling, and symmetry operations of the magnetic field.
It also offers to repeat the set-up periodically. In this example, we make
use of these options in order to
calculate the spontanous synchrotron radiation of a very long
undulator section including phase shifters. But first, we consider a
single cell of this section (Fig. \ref{ex7fig1}). We enter:

\cboxs{\senv \grau \\
python3 ../python/wave\_example\_7.py\\
cp \$WAVE/input/wave\_example\_7.in wave.in\\
wave\\
wplot \\
}

\begin{figure}[htb]
    \centering
    %\framebox(12,8){
    \includegraphics[width=0.8\textwidth]
    {/home/scheer/wav/doc/pp/wave_example_7_single_id.pdf}
    %}
    \caption{Field and 9th harmonic of a single undulator section}
    \label{ex7fig1}
\end{figure}

As the overview shows,
the field is a superposition a planar undulator (namelist \scap{\$ELLIPN}) and
a phase shifter (namelist \scap{\$HALBASY}). If \scap{IBSUPER=1}
in \wie, \W~accepts several field selections . The undulator parameters are
set such, that the 9th harmonic is at \gtis{150\,keV} for a beam energy of
\gtis{17\,GeV}. We now make use of a the periodic field option. Have a
look on the corresponding variables in the namelist
\scap{\$B0SCGLOBN}:\\

\cboxss{IPERIODG=1\\
PERIODG=3.\\
PEROFFG=-1.5\\
}

This means, that the magnetic field is periodically repeated starting from
\scap{X=-1.5} with a period-length of 3 meters.

We set in \wie:

\cboxss{XSTOP=88.5\\
IHTRSMP=10}

and run \W and get the results shown in Fig. \ref{ex7fig3}

\begin{figure}[htb]
    \centering
    %\framebox(12,8){
    \includegraphics[width=0.8\textwidth]
    {/home/scheer/wav/doc/pp/wave_example_7_long.pdf}
    %}
    \caption{Field and 9th harmonic of a long undulator section}
    \label{ex7fig3}
\end{figure}

Passing a long undulator section an electron beam loses a significant
amount of energy. Hence the K-parameter of the down-stream undulators
does not match the beam energy anymore. This results in a broadening
of the spectrum. This can be simulated by setting

\cboxss{IENELOSS=1}

in \wie. We run \W and get the results in Fig. \ref{ex7fig4}.

\begin{figure}[htb]
    \centering
    %\framebox(12,8){
    \includegraphics[width=0.8\textwidth]
    {/home/scheer/wav/doc/pp/wave_example_7_ieneloss.pdf}
    %}
    \caption{Field and 9th harmonic of a long undulator section
    with continous energy loss of the beam}
    \label{ex7fig4}
\end{figure}

To compensate it, we must adjust the field of each undulator
accordingly. This can be achieved with the field option

\cboxss{JBMASK=100}

and by providing a file \gtis{wave.bmask}, which contains a number of
$X$-Intervalls and a scaling factor for the magnetic field for
each intervall. The last \wo gives the total and relative energy loss of the beam due to
the 30 undulators. From this number we calculate the energy loss of a
single undulator cell, and use the relation

\begin{eqnarray*}
   \frac{E_{beam}^2}{1+K^2/2} &=& const
\end{eqnarray*}

to calculate the correction factor for the undulator field, as it is
done in the function \gtis{taper()}, which also writes the file \scap{wave.bmask}

We use this function:

\cboxss{wplot\\
taper()\\
Quit()
}

The default values of this function correspond to this example. Beside the
relative energy loss the starting $K$-value is taken from \wo of the
previous run as well (Look for ''deflection parameter'').

In \wi we set

\cboxss{JBMASK=100}

Please, also read the comments of the corresponding namelist
\scap{B0SCGLOBN} carefully and notice the difference between
\scap{JBMASK} and \scap{IBMASK}.

We run \W~again and see in Fig. \ref{ex7fig5} that the harmonic of the
undulator recovers.

\begin{figure}[htb]
    \centering
    %\framebox(12,8){
    \includegraphics[width=0.8\textwidth]
    {/home/scheer/wav/doc/pp/wave_example_7_taper.pdf}
    %}
    \caption{Field and 9th harmonic of a long undulator section
    with continous energy loss of the beam and corrected undulators}
    \label{ex7fig5}
\end{figure}
 % Global Field Options
\include{bibliography}
\end{document}
